% Notations du Tome III, communes aux deux langues (\EF). Ossature du 2026-10-29 : les conventions
% (celles des Tomes I et II) et le gabarit du tableau des noms. Les noms et les symboles se relèvent dans
% les chapitres une fois écrits (comme l'intégrateur du Tome II l'a fait) : une lettre qui sert deux fois
% est dite, avec la portée de chaque sens. Candidats déjà repérés à l'ossature, à définir par les
% rédacteurs là où ils servent, puis à porter ici : brin, codon, codon synonyme du zéro, fondu −0 → +0,
% mutation synonyme, S*, frontière (base frontière), λ (poids de la mutation synonyme), κ (classe de
% contexte ; et κ = σ/E[m], nombre de garde, au chapitre 11 : deux sens), H (hessienne des entrées
% d'un expert), e_k et λ_k (base propre de H ; λ y a un second sens), lots A et B, J1 et R1 (le jalon,
% les fenêtres de référence ; R1 est aussi la réserve R1 de l'inventaire des figures : trois sens),
% B2 (critère, et lemme B2 de la synthèse), P1 à P6 (pistes du mimétisme ; P1 et P2 sont aussi des
% figures de la piste des spectres), T1 à T7, F1 à F10, F-S1…, F-J1… (figures des notes, renvoyées aux
% numéros III-x.y), MXFP4, E2M1, E8M0, KL.
\chapter*{\EF{Notation and conventions}{Notations et conventions}}
\addcontentsline{toc}{chapter}{\EF{Notation and conventions}{Notations et conventions}}
\markboth{\EF{Notation and conventions}{Notations et conventions}}{\EF{Notation and conventions}{Notations et conventions}}

\EF{The conventions are those of Volumes~I and~II. Units are decimal: 1\,GB $=10^9$ bytes; binary multiples are written KiB, MiB, GiB. Every number carries its provenance: \mes{} measured here on Stone's machine with a committed tool and log, \calc{} computed from named numbers, \litt{} literature (publisher's number, not reproduced). A sentence that a later measurement contradicted is struck out where it stood (\barre{like this}), with what replaced it. A file of the repository cited as the source of a number is set as \fichier{mesures/genome/}; its path is relative to the root of the \texttt{troncons} repository. Figures are numbered III-$x$.$y$ in the inventory of the figures (\fichier{memoire/tome3/INVENTAIRE-figures.md}); a figure drawn outside the common style says so in its caption.}%
{Les conventions sont celles des Tomes~I et~II. Les unités sont décimales : 1\,Go $=10^9$ octets ; les multiples binaires s'écrivent Kio, Mio, Gio. Chaque nombre porte sa provenance : \mes{} mesuré ici, sur la machine de Stone, par un outil et un journal commités ; \calc{} calculé à partir de nombres nommés ; \litt{} littérature (chiffre d'éditeur, non reproduit). Une phrase qu'une mesure plus récente a démentie est barrée à sa place (\barre{ainsi}), avec ce qui l'a remplacée. Un fichier du dépôt cité comme source d'un nombre s'écrit \fichier{mesures/genome/} ; son chemin part de la racine du dépôt \texttt{troncons}. Les figures portent dans l'inventaire des figures un numéro III-$x$.$y$ (\fichier{memoire/tome3/INVENTAIRE-figures.md}) ; une figure tracée hors du style commun le dit dans sa légende.}

\section*{\EF{Names of the work}{Les noms du travail}}

% Gabarit du Tome II (memoire/tome2/commun/notations.tex) : une ligne par nom, \EF dans chaque cellule.
% Lignes relevées dans les chapitres écrits et jugés (grave G5 de memoire/tome3/JUGEMENT-redaction-fr.md,
% 2026-10-29) ; chaque ligne renvoie au texte qui la fonde.
\renewcommand{\arraystretch}{1.15}
\begin{longtable}{@{}p{0.19\linewidth}p{0.76\linewidth}@{}}
\toprule
\EF{Name}{Nom} & \EF{Meaning}{Signification} \\
\midrule
\endhead
troncons & \EF{the program and the project; Volume~II describes the program}{le programme et le projet ; le Tome~II décrit le programme} \\
\EF{control, break}{témoin, cassure} & \EF{the control of a measurement is a case whose answer is known in advance; its break forces the control's condition to false, and the instrument must then fail; every new control is broken once before it is trusted (\cref{t3ecl:sec:instruments})}{le témoin d'une mesure est un cas dont la réponse est connue d'avance ; sa cassure force sa condition à faux, et l'instrument doit alors tomber ; tout témoin neuf est cassé une fois avant qu'on s'y fie (\cref{t3ecl:sec:instruments})} \\
20b, 120b & \EF{gpt-oss-20b (24 layers of 32 experts: 768 experts) and gpt-oss-120b (36 layers of 128 experts: \num{4608} experts), the two models of the volume, their expert weights in MXFP4}{gpt-oss-20b (24 couches de 32 experts : 768 experts) et gpt-oss-120b (36 couches de 128 experts : \num{4608} experts), les deux modèles du tome, aux poids d'experts en MXFP4} \\
MXFP4, E2M1, E8M0 & \EF{the block format of the expert weights: 32 codes E2M1 (4 bits: sign, two exponent bits, one mantissa bit; sixteen codes for fifteen values, $\pm0$ being equal) and one shared scale E8M0 (one byte, an exponent alone): 17 bytes per block (\cref{t3que:ch})}{le format en blocs des poids d'experts : 32 codes E2M1 (4 bits : signe, deux bits d'exposant, un de mantisse ; seize codes pour quinze valeurs, $\pm0$ étant égaux) et une échelle commune E8M0 (un octet, un exposant seul) : 17 octets par bloc (\cref{t3que:ch})} \\
KL & \EF{the judge of the volume: Kullback--Leibler divergence of the next-token law of the modified model against the f64 reference, $\mathrm{KL}_t = \sum_i p_{t,i}\log(p_{t,i}/q_{t,i})$ in nats at position $t$ (only $\log p_{t,i} > -16$ is summed), averaged over the scored positions (\cref{t3ecl:sec:kl}, \cref{t3syn:def:kl})}{le juge du tome : divergence de Kullback-Leibler de la loi du jeton suivant du modèle modifié contre la référence f64, $\mathrm{KL}_t = \sum_i p_{t,i}\log(p_{t,i}/q_{t,i})$ en nats à la position $t$ (seuls les $\log p_{t,i} > -16$ sont sommés), moyennée sur les positions notées (\cref{t3ecl:sec:kl}, \cref{t3syn:def:kl})} \\
\EF{f64 reference}{référence f64} & \EF{\texttt{trc-reference}, the exact computation of gpt-oss in double precision (Volume~II, lot R); its logits are the truth of every KL}{\texttt{trc-reference}, le calcul exact de gpt-oss en double précision (Tome~II, lot R) ; ses logits sont la vérité de toute KL} \\
\EF{fork}{fourche} & \EF{a modified copy of the reference for a need of the round (capturing expert inputs, substituting a base)}{une copie modifiée de la référence pour un besoin de la ronde (capturer les entrées des experts, leur substituer une base)} \\
R1 & \EF{the ten evaluation windows of wiki.test of Volume~II (windows 0 to 9, 512 tokens each); the KL scores \num{2550} positions there, 255 per window. Not to be confused with \emph{the reserve R1}}{les dix fenêtres d'évaluation de wiki.test du Tome~II (fenêtres 0 à 9, de 512 jetons chacune) ; la KL y note \num{2550} positions, 255 par fenêtre. À ne pas confondre avec \emph{la réserve R1}} \\
\EF{calibration}{calibration} & \EF{windows 10 to 21 of wiki.test (\num{6132} positions), where the mechanisms were tuned; disjoint from R1}{les fenêtres 10 à 21 de wiki.test (\num{6132} positions), où les mécanismes ont été réglés ; disjointes de R1} \\
\EF{fresh windows}{fenêtres neuves} & \EF{the ten windows of wiki.test, touched neither by the calibration nor by R1, on which the judge of the construction re-did the results (\cref{t3jug:sec:fenetres})}{les dix fenêtres de wiki.test que ni la calibration ni R1 ne touchent, où le juge de la construction a refait les résultats (\cref{t3jug:sec:fenetres})} \\
P0, P1, P2, P2-10, P3 & \EF{the KL protocols of the round (\cref{t3syn:tab:protocoles}): P0 lossless (bit-identical logits); P1 the scout, 32 experts out of 768 on R1; P2 the genome, 8 experts on 3 windows; P2-10 the same 8 on the 10 windows of R1 (pass K10); P3 predictive coding and pruning, 128 experts on windows 5 to 9. A KL is compared across protocols only as a ratio to the base of its own protocol}{les protocoles de KL de la ronde (\cref{t3syn:tab:protocoles}) : P0 sans perte (logits identiques au bit) ; P1 l'éclaireur, 32 experts sur 768 sur R1 ; P2 le génome, 8 experts sur 3 fenêtres ; P2-10 les mêmes 8 sur les 10 fenêtres de R1 (passe K10) ; P3 le codage prédictif et l'élagage, 128 experts sur les fenêtres 5 à 9. Une KL ne se compare d'un protocole à l'autre qu'en rapport à la base de son propre protocole} \\
K10 & \EF{the pass that re-did $S^{*}$ against the base on the 10 windows of R1 (protocol P2-10), declared before; it replaces the 3-window ratio (reserve R4)}{la passe qui a refait $S^{*}$ contre la base sur les 10 fenêtres de R1 (protocole P2-10), déclarée avant ; elle remplace le rapport à 3 fenêtres (réserve R4)} \\
\EF{lots R, A, B}{lots R, A, B} & \EF{R: the reference and the judge of Volume~II; A: the zero-synonym codon in the format v1.2 (\cref{t3cod:ch}); B: the synonymous genome and the fast mode (\cref{t3sno:ch})}{R : la référence et le juge du Tome~II ; A : le codon synonyme du zéro dans le format v1.2 (\cref{t3cod:ch}) ; B : le génome synonyme et le mode rapide (\cref{t3sno:ch})} \\
\EF{arbiter, F1}{arbitre, F1} & \EF{the arbiter is the role of Volume~II that settled the format; F1 is its measurement of the frontier (rANS codes per lane, scales as offsets from the row median), reported in Volume~II}{l'arbitre est le rôle du Tome~II qui a tranché le format ; F1 est sa mesure de la frontière (codes en rANS par voie, échelles en écart à la médiane de ligne), rapportée au Tome~II} \\
\EF{frontier}{frontière} & \EF{the two points the round had to beat: lossless, 3.8431 bits per weight (A\_ctx) or 3.8434 (A\_flux); lossy, the \emph{base} of the fast mode (sign and four magnitude classes), 2.6441 bits read at 4.29\,\% relative weight error (\cref{t3que:sec:frontiere})}{les deux points que la ronde devait battre : sans perte, 3,8431 bits par poids (A\_ctx) ou 3,8434 (A\_flux) ; avec perte, la \emph{base} du mode rapide (signe et quatre classes de magnitude), 2,6441 bits lus à 4,29\,\% d'erreur relative sur les poids (\cref{t3que:sec:frontiere})} \\
A\_ctx, A\_flux & \EF{two accountings of the same lossless code: A\_ctx counts the codes coded by rANS per lane and per scale context, the scales, tables and headers; A\_flux counts the stream of the codec of study 03, as the arbiter accounts it. They differ by $3\cdot10^{-4}$ bit per weight}{deux comptabilités du même code sans perte : A\_ctx compte les codes en rANS par voie et par contexte d'échelle, les échelles, les tables et les en-têtes ; A\_flux compte le flux du codec de l'étude 03, comme l'arbitre le compte. Elles diffèrent de $3\cdot10^{-4}$ bit par poids} \\
\EF{strand, codon}{brin, codon} & \EF{in the genome of chapter~\ref{t3gen:ch}, the strand is the bit string of a row, read through a sliding window of $l$ bits (the codon) that a mixer and a table turn into a weight; elsewhere, \emph{codon} also names a 4-bit E2M1 code. Analogies, not arguments}{dans le génome du \cref{t3gen:ch}, le brin est la suite de bits d'une rangée, lue par une fenêtre glissante de $l$ bits (le codon) qu'un mélangeur et une table changent en poids ; ailleurs, \emph{codon} désigne aussi un code E2M1 de 4 bits. Des analogies, pas des arguments} \\
\EF{zero-synonym codon, fold}{codon synonyme du zéro, fondu} & \EF{codes 0 ($+0$) and 8 ($-0$) give the same value; the fold $-0 \to +0$ rewrites every 8 as 0 before coding, without changing a single value (\cref{t3ecl:sec:zero}, \cref{t3cod:ch})}{les codes 0 ($+0$) et 8 ($-0$) donnent la même valeur ; le fondu $-0 \to +0$ réécrit chaque 8 en 0 avant le codage, sans changer une seule valeur (\cref{t3ecl:sec:zero}, \cref{t3cod:ch})} \\
\EF{synonymous mutation}{mutation synonyme} & \EF{the choice, among the symbols of the base (same decoder), of those that minimise the output error weighted by $H$ rather than the weight error (\cref{t3ecl:sec:mutation})}{le choix, parmi les symboles de la base (même décodeur), de ceux qui minimisent l'erreur de sortie pondérée par $H$ plutôt que l'erreur sur les poids (\cref{t3ecl:sec:mutation})} \\
$H$, $e_k$, $\lambda_k$ & \EF{$H$ the (damped) second-moment matrix of the inputs $x$ of an expert matrix, weighted by the router weight $g$ ($\sum g^2 x x^\top / \sum g^2$), $e_k$ and $\lambda_k$ its eigenvectors and eigenvalues}{$H$ la matrice des moments d'ordre deux (amortie) des entrées $x$ d'une matrice d'expert, pondérées par le poids de routeur $g$ ($\sum g^2 x x^\top / \sum g^2$), $e_k$ et $\lambda_k$ ses vecteurs et valeurs propres} \\
$\lambda$ & \EF{the weight of the entropy constraint of the synonymous mutation ($\lambda = 0$: none; $\lambda = 1$); not to be confused with $\lambda_k$}{le poids de la contrainte d'entropie de la mutation synonyme ($\lambda = 0$ : aucune ; $\lambda = 1$) ; à ne pas confondre avec $\lambda_k$} \\
$\delta$, $\delta^{*}$ & \EF{the damping added to $H$; $\delta^{*} = 1$ the value chosen without looking at the evaluation; in \cref{t3syn:ch}, $\delta$ also writes the scale offset that sets $\kappa$}{l'amortissement ajouté à $H$ ; $\delta^{*} = 1$ la valeur choisie sans regarder l'évaluation ; au \cref{t3syn:ch}, $\delta$ note aussi l'écart d'échelle qui fixe $\kappa$} \\
$\kappa$ & \EF{two meanings: the scale class of a block, three classes set by its E8M0 against the row median (chapters~\ref{t3ecl:ch} to~\ref{t3sno:ch}); in \cref{t3tro:ch}, the guard number $\sigma/\E[m]$ of the router}{deux sens : la classe d'échelle d'un bloc, trois classes selon son E8M0 rapporté à la médiane de la ligne (\crefrange{t3ecl:ch}{t3sno:ch}) ; au \cref{t3tro:ch}, le nombre de garde $\sigma/\E[m]$ du routeur} \\
$S^{*}$ & \EF{the system of the genome: strand, rotation, weighted Viterbi and three exact exceptions per row (\cref{t3gen:ch})}{le système du génome : brin, rotation, Viterbi pondéré et trois exceptions exactes par rangée (\cref{t3gen:ch})} \\
v1.0, v1.2 & \EF{versions of the \texttt{.trc} format: v1.0 the original files; v1.2 adds the fold of the zero (\cref{t3cod:sec:format}); v1.1 (two levels) is reserved and written by no coder}{versions du format \texttt{.trc} : v1.0 les fichiers d'origine ; v1.2 y ajoute le fondu du zéro (\cref{t3cod:sec:format}) ; la v1.1 (deux niveaux) est réservée et n'est écrite par aucun codeur} \\
\EF{fast companion}{compagnon rapide} & \EF{the file \texttt{<name>.rapide.trc} that the fast mode of the engine reads, a plain \texttt{.trc} holding the synonym base (\cref{t3sno:sec:format})}{le fichier \texttt{<nom>.rapide.trc} que lit le mode rapide du moteur, un \texttt{.trc} ordinaire qui porte la base synonyme (\cref{t3sno:sec:format})} \\
J1 & \EF{the milestone of Volume~II, the exact service of gpt-oss-120b in a bounded window; here, its generation speed re-measured (\cref{t3cod:sec:j1})}{le jalon du Tome~II, le service exact de gpt-oss-120b dans une fenêtre bornée ; ici, sa vitesse de génération refaite (\cref{t3cod:sec:j1})} \\
ABBA, $r(B/A)$ & \EF{a timed series alternating the two files (A B B A A B B A), each pass a fresh process; $r(B/A)$ the ratio of their median speeds, judged against $2\sigma$ of the series}{une série chronométrée qui alterne les deux fichiers (A B B A A B B A), chaque passe dans un processus neuf ; $r(B/A)$ le rapport de leurs vitesses médianes, jugé contre $2\sigma$ de la série} \\
\EF{reserves R1 to R6}{réserves R1 à R6} & \EF{the conditions of the inventory of the figures that keep a figure or a sentence true (\cref{t3int:sec:lire}); R1: the $\times 1.44$ of the J1 comes from the format v1.2 through the pattern of its read requests, on this disk}{les conditions de l'inventaire des figures qui gardent vraie une figure ou une phrase (\cref{t3int:sec:lire}) ; R1 : le $\times 1{,}44$ du J1 vient du format v1.2 par le motif de ses requêtes de lecture, sur ce disque} \\
III-$x$.$y$ & \EF{the number of a figure in the inventory (\fichier{memoire/tome3/INVENTAIRE-figures.md}); the printed number follows the chapters of the PDF, where the introduction is chapter~1}{le numéro d'une figure dans l'inventaire (\fichier{memoire/tome3/INVENTAIRE-figures.md}) ; le numéro imprimé suit les chapitres du PDF, où l'introduction est le chapitre~1} \\
\bottomrule
\end{longtable}
